\chapter{Operators, Measurement, and Wave Functions}

Quantum mechanics is linear algebra. A state is a vector, while a measurable
quantity is represented by a special kind of linear operator. The purpose of
this lecture is to connect that abstract statement to calculations. We shall
first turn operators into matrices, then identify Hermitian operators as
observables, state the measurement postulates, and finally apply the entire
structure to a particle moving on a line.

\section{Operators, Matrix Elements, and Bases}

A linear operator maps one ket into another:
\begin{equation}
  \hat K\ket{A}=\ket{C}.
\end{equation}
Once a bra is supplied, the result becomes a complex number,
\begin{equation}
  \bra{B}\hat K\ket{A}.
\end{equation}
The order is important. The operator first acts on \(\ket A\); the resulting
ket is then paired with \(\bra B\).

Choose an orthonormal basis \(\{\ket n\}_{n=1}^{d}\). An ordinary
three-dimensional sketch helps us remember what a basis does, although a
quantum state space is generally complex and may have far more than three
dimensions.

\begin{figure}[tbp]
  \centering
  \includegraphics[width=0.46\textwidth]{lecture_03_figure_02.png}
  \caption{An ordinary vector sketch standing in for an orthonormal basis.}
\end{figure}
\lectureframe{00:04:29}

\begin{figure}[tbp]
  \centering
  \begin{tikzpicture}[scale=1.1,>=Latex]
    \draw[->,thick] (0,0) -- (0,1.55);
    \draw[->,thick] (0,0) -- (1.65,0);
    \draw[->,thick] (0,0) -- (-1.05,-1.05);
    \fill (0,0) circle (1.2pt);
  \end{tikzpicture}
  \caption{A clean reconstruction of the three-direction basis sketch.}
\end{figure}

There are many possible bases. A vector drawn in one basis can always be
expanded in another complete basis; only its components change.

\begin{classroomqa}
\audiencequestion{Can a vector associated with one drawn basis be represented
in the other, and do the labels \(m\) and \(n\) refer to different
bases?\lecturetimestamp{00:05:24--00:07:40}}
\lecturerresponse{Yes, the same vector can be expanded in either complete
basis. In the matrix element below, however, both \(m\) and \(n\) are labels
drawn from one chosen basis. They identify the row and column of the matrix
representing \(\hat K\) in that basis.}
\end{classroomqa}

The basis matrix elements are
\begin{equation}
  K_{mn}:=\bra m\hat K\ket n.
\end{equation}
There are \(d^2\) of them in a \(d\)-dimensional space, and together they
determine the operator once the basis has been fixed.

\begin{figure}[tbp]
  \centering
  \includegraphics[width=0.48\textwidth]{lecture_03_figure_03.png}
  \caption{The matrix element \(\bra m\hat K\ket n\) on the board.}
\end{figure}
\lectureframe{00:06:22}

\section{From Kets to Matrix Multiplication}

Every ket can be expanded in the chosen basis:
\begin{equation}
  \ket A=\sum_n A_n\ket n,
  \qquad
  A_n=\braket nA.
\end{equation}
Substitution gives
\begin{equation}
  \ket A=\sum_n\ket n\bra n\ket A.
\end{equation}
Since this holds for every \(\ket A\), the operator between the summation sign
and the final ket must be the identity:
\begin{equation}
  \sum_n\ket n\bra n=\mathbb I.
  \label{eq:discrete-completeness}
\end{equation}
This is the completeness relation. The adjacent ket and bra form a projector
onto one basis direction; summing all such projectors does nothing to the
vector.

We can now calculate the components of \(\hat K\ket A\):
\begin{align}
  \bra n\hat K\ket A
  &=\sum_m\bra n\hat K\ket m\braket mA \\
  &=\sum_m K_{nm}A_m.
  \label{eq:matrix-action}
\end{align}
Thus the abstract operation \(\hat K\ket A\) is ordinary matrix
multiplication when written in components.

\begin{figure}[tbp]
  \centering
  \includegraphics[width=0.72\textwidth]{lecture_03_figure_04.png}
  \caption{The passage from operator action to the component sum
  \(\sum_m K_{nm}A_m\).}
\end{figure}
\lectureframe{00:13:37}

The same correspondence holds for products of operators. In \(\hat K\hat L\),
\(\hat L\) acts first and \(\hat K\) acts second. Inserting
Eq.~\eqref{eq:discrete-completeness} between them gives
\begin{align}
  (KL)_{nm}
  &=\bra n\hat K\hat L\ket m \\
  &=\sum_r\bra n\hat K\ket r\bra r\hat L\ket m \\
  &=\sum_r K_{nr}L_{rm}.
  \label{eq:operator-product}
\end{align}
This is precisely the rule for multiplying matrices. It also makes clear why
operator order matters: in general \(\hat K\hat L\neq\hat L\hat K\).

\begin{figure}[tbp]
  \centering
  \includegraphics[width=0.72\textwidth]{lecture_03_figure_05.png}
  \caption{Insertion of a complete basis between two operators.}
\end{figure}
\lectureframe{00:18:30}

\section{Hermitian Operators and Their Eigenvectors}

The operators associated with observables must have real measurement values.
The relevant class is the Hermitian operators. Here the letter \(H\) merely
stands for Hermitian; it is not yet the Hamiltonian.

\begin{definition}
An operator \(\hat H\) is Hermitian when
\begin{equation}
  \bra B\hat H\ket A
  =
  \left(\bra A\hat H\ket B\right)^*
  \label{eq:hermitian-definition}
\end{equation}
for vectors in its domain. In finite dimensions this is equivalent to saying
that \(\bra A\hat H\ket A\) is real for every \(\ket A\).
\end{definition}

Most vectors change direction under a general operator. An eigenvector is a
special direction that does not:
\begin{equation}
  \hat H\ket\lambda=\lambda\ket\lambda.
  \label{eq:eigenvalue}
\end{equation}
Only its length and phase may change, through the scalar eigenvalue
\(\lambda\).

\begin{classroomqa}
\audiencequestion{Does every operation have eigenvectors?
\lecturetimestamp{00:29:53--00:30:18}}
\lecturerresponse{Over a finite-dimensional complex vector space, every
linear operator has at least one eigenvector, but a general operator need not
have enough eigenvectors to form a basis, much less an orthonormal basis. A
real rotation can also lack real eigenvectors. The special result needed here
is stronger: a Hermitian operator has a complete orthonormal eigenbasis. In
infinite-dimensional spaces, continuous spectra require generalized
eigenvectors, as the position example will show.}
\end{classroomqa}

The finite-dimensional spectral theorem supplies three facts:
\begin{enumerate}
  \item every eigenvalue of a Hermitian operator is real;
  \item eigenvectors with distinct eigenvalues are orthogonal;
  \item the space has a complete orthonormal basis of eigenvectors.
\end{enumerate}

The first two statements follow directly from Hermiticity. If
\(\hat H\ket\lambda=\lambda\ket\lambda\), then
\begin{equation}
  \bra\lambda\hat H\ket\lambda
  =\lambda\braket\lambda\lambda.
\end{equation}
The left side is real and the norm is positive, so \(\lambda\) is real.
For two eigenvectors,
\begin{align}
  \bra{\lambda_1}\hat H\ket{\lambda_2}
  &=\lambda_2\braket{\lambda_1}{\lambda_2},\\
  \bra{\lambda_1}\hat H\ket{\lambda_2}
  &=\lambda_1\braket{\lambda_1}{\lambda_2}.
\end{align}
Consequently,
\begin{equation}
  (\lambda_2-\lambda_1)\braket{\lambda_1}{\lambda_2}=0.
\end{equation}
Distinct eigenvalues therefore give orthogonal eigenvectors. If an eigenvalue
is degenerate, an orthonormal basis can be chosen within its eigenspace.
Finally, every eigenvector can be rescaled to unit norm.

\begin{figure}[tbp]
  \centering
  \includegraphics[width=0.62\textwidth]{lecture_03_figure_07.png}
  \caption{The eigenvalue equation and the three spectral facts for a
  Hermitian operator.}
\end{figure}
\lectureframe{00:33:40}

\begin{classroomqa}
\audiencequestion{Is a rotation Hermitian?
\lecturetimestamp{00:35:35--00:36:55}}
\lecturerresponse{A finite rotation is generally unitary, or orthogonal in a
real vector space, rather than Hermitian. Its infinitesimal generator is
anti-Hermitian; equivalently, the generator can be written as \(-i\) times a
Hermitian operator. In a real basis, symmetric matrices are the analogues of
Hermitian matrices, while infinitesimal rotation generators are
antisymmetric.}
\end{classroomqa}

\begin{classroomqa}
\audiencequestion{In which vector space do the eigenvectors form a basis?
\lecturetimestamp{00:37:35--00:40:05}}
\lecturerresponse{They form a basis of the same complex vector space on which
\(\hat H\) acts. If we choose that eigenbasis for our coordinates, the matrix
of \(\hat H\) is diagonal, with its eigenvalues along the diagonal.}
\end{classroomqa}

\section{Five Measurement Postulates}

The basic rules can now be stated in the same language.

\begin{enumerate}
  \item A physical state is represented by a ket in a complex vector space.
  A coin has a two-dimensional state space, a die has a six-dimensional one,
  and a particle on a line requires an infinite-dimensional space.

  \item Every observable is represented by a Hermitian operator.

  \item The possible outcomes of measuring an observable are the eigenvalues
  of its operator.

  \item An eigenvector \(\ket{\lambda_n}\) is a state in which that observable
  has the definite value \(\lambda_n\).

  \item If the normalized state is \(\ket A\), the probability of obtaining
  \(\lambda_n\) is
  \begin{equation}
    P(\lambda_n\mid A)
    =\left|\braket{\lambda_n}A\right|^2.
    \label{eq:born-discrete}
  \end{equation}
\end{enumerate}

\begin{figure}[tbp]
  \centering
  \includegraphics[width=0.58\textwidth]{lecture_03_figure_08.png}
  \caption{The fifth postulate written as a squared projection onto an
  eigenvector.}
\end{figure}
\lectureframe{00:54:50}

If the eigenvectors form an orthonormal basis, normalization guarantees that
the probabilities sum to one:
\begin{equation}
  \sum_n P(\lambda_n\mid A)
  =\sum_n\left|\braket{\lambda_n}A\right|^2
  =\braket AA=1.
\end{equation}

\begin{classroomqa}
\audiencequestion{Is a state itself an observable?
\lecturetimestamp{00:42:05--00:42:50}}
\lecturerresponse{No. Position, momentum, spin components, and fields are
examples of observables. A state is not a number read from an apparatus; it
encodes the probabilities for the possible readings. A single measurement
therefore does not reveal the complete state.}
\end{classroomqa}

\begin{classroomqa}
\audiencequestion{If eigenvalues are real, what about experimental error
bars?\lecturetimestamp{00:45:33--00:46:28}}
\lecturerresponse{Real eigenvalues do not eliminate experimental uncertainty.
An apparatus may report a real interval or a real value with an error bar.
The point is that a measurement outcome is not an imaginary number.}
\end{classroomqa}

\begin{classroomqa}
\audiencequestion{Do we observe eigenvectors, and what does it mean for an
observable to be definite?\lecturetimestamp{00:47:41--00:50:29}}
\lecturerresponse{The apparatus records an eigenvalue, not an eigenvector. An
observable is definite when repeated ideal measurements on identically
prepared systems give the same value with probability one. For example, a
spin state can be definite along one axis and probabilistic along another,
because it is an eigenstate of one spin operator but not of the other.}
\end{classroomqa}

\section{A Particle on a Line}

We now leave finite-dimensional examples and consider a particle moving on
the real line. Its states are represented by complex functions \(\psi(x)\).
Functions form a vector space, and the natural inner product is
\begin{equation}
  \braket\phi\psi
  =
  \int_{-\infty}^{\infty}\phi^*(x)\psi(x)\,dx.
  \label{eq:function-inner-product}
\end{equation}
We shall initially work with square-integrable, sufficiently well-behaved
functions and impose whatever boundary conditions are needed by the
operators.

\begin{classroomqa}
\audiencequestion{Is the inner product of two wave functions always real?
\lecturetimestamp{00:58:13--00:59:07}}
\lecturerresponse{No. The inner product \(\braket\phi\psi\) can be complex.
What must be real and nonnegative is a vector's inner product with itself:
\[
  \braket\psi\psi=\int_{-\infty}^{\infty}|\psi(x)|^2\,dx.
\]}
\end{classroomqa}

\subsection{The position operator}

The position operator multiplies a wave function by its argument:
\begin{equation}
  (\hat x\psi)(x)=x\psi(x).
  \label{eq:position-operator}
\end{equation}
It is Hermitian on its proper domain because
\begin{equation}
  \bra\psi\hat x\ket\psi
  =
  \int_{-\infty}^{\infty}x|\psi(x)|^2\,dx
\end{equation}
is real.

The position eigenvalue equation is
\begin{equation}
  \hat x\ket\lambda=\lambda\ket\lambda,
\end{equation}
or, in the \(x\)-representation,
\begin{equation}
  (x-\lambda)\psi_\lambda(x)=0.
  \label{eq:position-eigenvalue}
\end{equation}
For every \(x\neq\lambda\), the factor \(x-\lambda\) is nonzero, so the
eigenfunction must vanish. Its support is concentrated at the single point
\(x=\lambda\).

\begin{classroomqa}
\audiencequestion{Have we assumed that the norm integral converges?
\lecturetimestamp{01:05:02--01:05:42}}
\lecturerresponse{Yes. Physical state vectors are normalized, so their
wavefunctions are square-integrable. For the elementary manipulations here
we also assume sufficiently tame behavior at large distance. Square
integrability by itself does not force pointwise decay in every mathematical
sense, but the chosen operator domains will supply the boundary behavior
needed below.}
\end{classroomqa}

No ordinary square-integrable function can be nonzero at only one point and
have unit norm. The position eigenfunction is therefore a generalized
function:
\begin{equation}
  \braket{x}{\lambda}=\delta(x-\lambda).
  \label{eq:position-delta}
\end{equation}
The lecture pictures \(\delta(x-\lambda)\) as a rectangle of width
\(\epsilon\) and height \(1/\epsilon\), then lets \(\epsilon\) shrink. This
choice gives unit area,
\begin{equation}
  \int_{-\infty}^{\infty}\delta(x-\lambda)\,dx=1,
\end{equation}
but it does not give unit Hilbert-space norm: the integral of the square of
the rectangle grows like \(1/\epsilon\). The delta function is not an
ordinary normalized state.

\begin{figure}[tbp]
  \centering
  \includegraphics[width=0.62\textwidth]{lecture_03_figure_09.png}
  \caption{A narrow unit-area approximation to
  \(\delta(x-\lambda)\).}
\end{figure}
\lectureframe{01:08:33}

\begin{classroomqa}
\audiencequestion{Are these position eigenfunctions orthonormal?
\lecturetimestamp{01:09:13--01:11:22}}
\lecturerresponse{They are orthogonal in the generalized continuous sense,
but not normalizable one by one as ordinary vectors. The precise replacement
for Kronecker-delta normalization is
\[
  \braket{x}{x'}=\delta(x-x').
\]
Position eigenkets are delta-normalized generalized eigenvectors.}
\end{classroomqa}

There is one such generalized eigenvector for every real position, so the
spectrum of \(\hat x\) is continuous. Completeness now takes the form
\begin{equation}
  \int_{-\infty}^{\infty}\ket x\bra x\,dx=\mathbb I.
  \label{eq:continuous-completeness}
\end{equation}

\subsection{What the wave function means}

The delta function samples an ordinary function:
\begin{equation}
  \int_{-\infty}^{\infty}
  \delta(x-\lambda)f(x)\,dx=f(\lambda).
  \label{eq:delta-sampling}
\end{equation}
Applying this to a state gives
\begin{equation}
  \braket\lambda\psi
  =
  \int_{-\infty}^{\infty}
  \delta(x-\lambda)\psi(x)\,dx
  =
  \psi(\lambda).
\end{equation}
Relabeling the eigenvalue by \(x\),
\begin{equation}
  \psi(x)=\braket x\psi.
  \label{eq:wavefunction-component}
\end{equation}
The wave function is therefore the component of the abstract state ket in
the position basis.

The fifth postulate then says that
\begin{equation}
  \rho(x)=|\psi(x)|^2
  \label{eq:position-density}
\end{equation}
is the position probability density. A probability belongs to an interval,
not to one isolated point:
\begin{equation}
  P(a\leq x\leq b)
  =
  \int_a^b|\psi(x)|^2\,dx.
\end{equation}

\begin{classroomqa}
\audiencequestion{Has \(\psi\) changed from an arbitrary function into one
special function?\lecturetimestamp{01:18:43--01:19:28}}
\lecturerresponse{No. Any normalized state wavefunction in the allowed
domain supplies a probability distribution through \(|\psi(x)|^2\).
The wavefunction itself may be complex and may carry signs and phases. Its
absolute square is real and nonnegative. Those phases are essential because
position is not the only observable we may measure; replacing \(\psi\) by
the probability density alone would discard information needed for other
measurements.}
\end{classroomqa}

\section{Differentiation and the Momentum Operator}

Differentiation is another linear operation on functions:
\begin{equation}
  \frac{d}{dx}(\alpha\psi+\beta\phi)
  =
  \alpha\frac{d\psi}{dx}
  +
  \beta\frac{d\phi}{dx}.
\end{equation}
The derivative itself is not Hermitian. With boundary terms chosen to vanish,
it is anti-Hermitian:
\begin{equation}
  \left(\frac{d}{dx}\right)^\dagger
  =
  -\frac{d}{dx}.
\end{equation}
Multiplication by \(-i\) converts it into a Hermitian operator,
\begin{equation}
  \hat K=-i\frac{d}{dx}.
  \label{eq:k-operator}
\end{equation}

The key calculation is integration by parts. For functions in a common
domain with \(\phi^*(x)\psi(x)\to0\) at both ends,
\begin{align}
  \bra\phi\hat K\ket\psi
  &=
  -i\int_{-\infty}^{\infty}
  \phi^*(x)\frac{d\psi}{dx}\,dx\\
  &=
  -i\left[\phi^*(x)\psi(x)\right]_{-\infty}^{\infty}
  +i\int_{-\infty}^{\infty}
  \frac{d\phi^*}{dx}\psi(x)\,dx\\
  &=
  \braket{\hat K\phi}{\psi}.
  \label{eq:k-hermitian}
\end{align}
Thus \(\hat K\) is Hermitian on that domain. The boundary conditions are part
of the operator; without them, the formal differential expression alone is
not enough to establish Hermiticity.

\begin{figure}[tbp]
  \centering
  \includegraphics[width=0.68\textwidth]{lecture_03_figure_10.png}
  \caption{The integration-by-parts step that moves the derivative and
  changes its sign.}
\end{figure}
\lectureframe{01:31:37}

\subsection{Plane-wave eigenfunctions}

The eigenvalue equation
\begin{equation}
  -i\frac{d\psi_k}{dx}=k\psi_k
\end{equation}
has the solution
\begin{equation}
  \psi_k(x)\propto e^{ikx}
  =\cos(kx)+i\sin(kx).
  \label{eq:plane-wave}
\end{equation}
Indeed,
\begin{equation}
  -i\frac{d}{dx}e^{ikx}
  =
  -i(ik)e^{ikx}
  =
  k e^{ikx}.
\end{equation}

\begin{figure}[tbp]
  \centering
  \includegraphics[width=0.68\textwidth]{lecture_03_figure_11.png}
  \caption{The plane-wave eigenfunction of
  \(\hat K=-i\,d/dx\).}
\end{figure}
\lectureframe{01:34:50}

A plane wave is the opposite of a position eigenstate. Its magnitude is
constant:
\begin{equation}
  |e^{ikx}|^2=1.
\end{equation}
Formally, it gives no preferred position at all. It is also not
square-integrable on the whole line, so, like the delta-localized position
eigenket, it is a generalized eigenstate rather than a normalizable physical
state by itself.

\begin{classroomqa}
\audiencequestion{Did the Hermiticity argument not assume decay at infinity,
while the plane waves fail to decay?\lecturetimestamp{01:38:41--01:39:47}}
\lecturerresponse{Yes. The tension is real and should not be hidden. Plane
waves belong to the generalized continuous spectrum. One may approximate
them by long, normalizable wave packets, work in a large periodic box and
then take its size to infinity, or use delta normalization. The formal plane
waves remain useful because they expose the spectral structure exactly.}
\end{classroomqa}

\section{Why Wave Number Measures Momentum}

One full oscillation of \(e^{ikx}\) occurs when its phase changes by \(2\pi\).
If \(L\) denotes the wavelength, then
\begin{equation}
  kL=2\pi,
  \qquad
  L=\frac{2\pi}{|k|}.
  \label{eq:wavelength}
\end{equation}
The sign of \(k\) distinguishes the two propagation directions, while its
magnitude fixes the wavelength.

Einstein's photon relation and de Broglie's extension to material particles
suggest
\begin{equation}
  p=\frac{h}{L}.
\end{equation}
Using \(L=2\pi/|k|\) for the magnitude and retaining the sign for direction,
\begin{equation}
  p=\hbar k,
  \qquad
  \hat p=\hbar\hat K=-i\hbar\frac{d}{dx}.
  \label{eq:momentum-operator}
\end{equation}

\begin{figure}[tbp]
  \centering
  \includegraphics[width=0.62\textwidth]{lecture_03_figure_12.png}
  \caption{Wavelength and the de Broglie identification
  \(p=\hbar k\).}
\end{figure}
\lectureframe{01:43:50}

At this stage, Eq.~\eqref{eq:momentum-operator} is a well-motivated
identification rather than a completed derivation from classical mechanics.
To justify the name physically, we must introduce time evolution and show
that localized wave packets move in the classical limit with velocity related
to their \(k\)-content.

\begin{classroomqa}
\audiencequestion{Does large \(k\) mean a large or a small wavelength?
\lecturetimestamp{01:54:03--01:54:30}}
\lecturerresponse{A large \(|k|\) gives rapid variation and a short
wavelength; a small \(|k|\) gives slow variation and a long wavelength,
because \(L=2\pi/|k|\).}
\end{classroomqa}

\section{Position, Momentum, and Incompatible Bases}

The position and momentum eigenfunctions could hardly look more different.
A position eigenfunction is concentrated at one point and contains no
definite \(k\). A momentum eigenfunction extends over the entire line and
contains no definite position. The two families form different generalized
bases in the same state space.

This is the first glimpse of the uncertainty principle in the language of
state vectors. A sharply localized wave packet requires a superposition of
many values of \(k\), while a single value of \(k\) produces complete
delocalization. The precise inequality will come later; the clash of the two
eigenbases is already visible.

There is also a useful geometric way to picture the complex plane wave. At
each point \(x\), \(e^{ikx}\) is a point on the unit circle in the complex
plane. As \(x\) increases, that point winds around the circle. If the spatial
axis is drawn perpendicular to the complex plane, the resulting curve is a
helix. Large \(|k|\) winds rapidly; small \(|k|\) winds slowly.

\begin{classroomqa}
\audiencequestion{Can \(e^{ikx}\) be visualized as a function twisting around
the \(x\)-axis?\lecturetimestamp{01:50:04--01:53:05}}
\lecturerresponse{Yes. The real and imaginary parts are the two perpendicular
coordinates of a unit vector in the complex plane. Advancing along the
spatial axis rotates that vector through phase \(kx\), producing the helix
described above.}
\end{classroomqa}

\begin{figure}[tbp]
  \centering
  \begin{tikzpicture}[>=Latex]
    \draw[->,gray!75] (0,0) -- (7.8,0) node[below right] {\(x\)};
    \draw[->,gray!75] (0,-0.9) -- (0,1.0)
      node[above] {complex amplitude};
    \draw[blue!65!black,thick,decorate,
      decoration={coil,aspect=0.72,segment length=8mm,amplitude=5mm}]
      (0.45,0) -- (7.15,0);
    \draw[->,gray!65] (0.45,0) -- (0.45,0.78)
      node[left] {\(\psi(0)\)};
  \end{tikzpicture}
  \caption{Transcript-based schematic of the plane-wave phase winding as a
  helix along the spatial axis.}
\end{figure}

\begin{classroomqa}
\audiencequestion{Why was wavefunction collapse absent from the five
postulates, and what is the probability distribution for momentum?
\lecturetimestamp{01:53:15--01:54:02}}
\lecturerresponse{The lecture has so far stated only which outcomes are
possible and how their probabilities are assigned. The physical measurement
process and its effect on the state have not yet been introduced. Likewise,
the momentum-space probability rule requires expanding the state in the
\(\ket k\) basis; that construction is deferred to the next step of the
course.}
\end{classroomqa}

\begin{classroomqa}
\audiencequestion{Classical momentum involves motion in time. Why has a
spatial derivative appeared before any time derivative?
\lecturetimestamp{01:54:31--01:56:03}}
\lecturerresponse{The identification must ultimately be checked through time
evolution. A normalizable wave packet forms a localized lump. In the
appropriate heavy-particle or classical limit, that lump should move
according to the classical equations, with its velocity related to the
wave-number content of the packet. Time has not yet entered the theory, so
that dynamical check remains ahead.}
\end{classroomqa}
